\chapter{แนวคิด ทฤษฎี และเอกสารงานวิจัยที่เกี่ยวข้อง}

ในการวิจัยเรื่อง ``การพัฒนาทักษะเกษตรกรอัจฉริยะและการจำแนกกลุ่มเป้าหมายด้วยเทคโนโลยีการเรียนรู้ของเครื่อง'' ผู้วิจัยได้ศึกษาแนวคิด ทฤษฎี และงานวิจัยที่เกี่ยวข้อง เพื่อเป็นพื้นฐานในการกำหนดกรอบแนวคิดการวิจัย โดยแบ่งหัวข้อการนำเสนอดังนี้

\begin{itemize}
\item บริบทภาคการเกษตรของไทยและนโยบายเกษตรกรอัจฉริยะ (Smart Farmer)
\item ความเป็นมาและวัตถุประสงค์ของนโยบาย Smart Farmer
\item ปัจจัยที่ส่งผลต่อความสามารถและการยอมรับเทคโนโลยีของเกษตรกร
\item บทบาทของ PSM และ DID ในการประเมินนโยบายเชิงเศรษฐศาสตร์
\item เทคโนโลยีการเรียนรู้ของเครื่อง (Machine Learning)
\end{itemize}

\section{บริบทและนโยบายเกษตรกรอัจฉริยะ (Smart Farmer)}

ภาคการเกษตรมีบทบาทสำคัญในการขับเคลื่อนเศรษฐกิจไทย โดยมีแรงงานมากกว่าหนึ่งในสามหรือประมาณร้อยละ 30 ของกำลังแรงงานทั้งหมดอยู่ในภาคการเกษตร \cite{un2563}

การจ้างงานในภาคเกษตรกรรมของไทยครอบคลุมถึง 6.4 ล้านครัวเรือน อย่างไรก็ตาม ภาคการเกษตรมีสัดส่วนเพียงร้อยละ 10 ของผลิตภัณฑ์มวลรวมในประเทศ (ปี พ.ศ. 2562) และมีอัตราการเติบโตที่ลดลงในช่วงสามทศวรรษที่ผ่านมา โดยภาคการเกษตรต้องเผชิญกับความท้าทายเชิงโครงสร้างและปัจจัยภายนอกที่สำคัญ ประกอบด้วย

\begin{itemize}

\item กับดักรายได้ต่ำและความเหลื่อมล้ำ: แม้ภาคเกษตรจะใช้แรงงานจำนวนมาก แต่กลับสร้างมูลค่าผลิตภัณฑ์มวลรวมในประเทศ (GDP) ได้เพียงร้อยละ 10 และเกษตรกรส่วนใหญ่ยังติดอยู่ในวงจรหนี้สินและการเป็นเพียงผู้ผลิตวัตถุดิบขั้นต้น (Raw Material) ที่ไม่สามารถกำหนดราคาเองได้ ส่งผลให้ขาดอำนาจต่อรองด้านราคาซื้อขาย \cite{meechoovet2023}

\item สังคมเกษตรสูงวัย (Aging Society): อายุเฉลี่ยของหัวหน้าครัวเรือนเกษตรกรไทยสูงถึง 62 ปี โดยเกษตรกลุ่มนี้มักมีแนวโน้มที่จะรักษาวิธีการผลิตแบบตั้งเดิมเอาไว้ เนื่องจากความสามารถยอมรับเทคโนโลยีใหม่ทำได้ช้าและมีความพร้อมในการลงทุนทางเทคโนโลยีต่ำ \cite{duang2019} ขณะที่แรงงานรุ่นใหม่ (Young Smart Farmer) ยังเข้าสู่ภาคเกษตรไม่เพียงพอที่จะทดแทนแรงงานเดิมที่ทยอยออกจากระบบ

\item ความผันผวนของสภาพภูมิอากาศและต้นทุน: การเปลี่ยนแปลงสภาพภูมิอากาศ (Climate Change) ก่อให้เกิดภัยแล้งและอุทกภัยที่รุนแรงและคาดการณ์ยาก ประกอบกับต้นทุนปัจจัยการผลิตที่พุ่งสูงขึ้น เช่น ราคาปุ๋ยเคมีและพลังงานในช่วงปี 2565-2566 กดดันให้รายได้สุทธิของเกษตรกรลดลง ด้วยเหตุผลดังกล่าว ส่งผลให้กลุ่มเกษตรที่มีความพร้อมทางความรู้และเทคโนโลยีดำเนินการปรับตัวไปสู่การทำเกษตรแบบแม่นยำ (Precision Agriculture) เพื่อลดความเสี่ยง \cite{hwang2024}

\end{itemize}

การชะลอตัวของการเติบโตดังกล่าวมีสาเหตุหลายประการ การลงทุนในการผลิตทางการเกษตรเป็นการลงทุนระยะยาวที่มีความเสี่ยงหลายด้าน เกษตรกรต้องเผชิญกับความไม่แน่นอนทางธรรมชาติ เช่น ภัยแล้งและอุทกภัย ปัญหาด้านทรัพยากร เช่น การขาดแคลนน้ำ มลพิษ และการเสื่อมโทรมของดิน รวมถึงความผันผวนของตลาดทั้งในด้านผลผลิตและราคา \cite{chadarat2562,un2563}

งานวิจัยหลายชิ้นพบว่า การนำเทคโนโลยีมาประยุกต์ใช้ช่วยให้เกษตรกรสามารถพัฒนาผลิตภาพและประสิทธิภาพในภาคการเกษตร อันนำไปสู่การเพิ่มรายได้จากการทำเกษตรกรรม \cite{jinjarak2550,rueatee2557,thongpan2556}

อย่างไรก็ตาม เกษตรกรไทยยังเผชิญความท้าทายหลายประการในการนำเทคโนโลยีมาใช้ โดยเฉพาะการเข้าถึงเทคโนโลยีที่จำกัด การขาดแคลนองค์ความรู้ทางเทคนิค และการสนับสนุนจากภาครัฐที่ไม่เพียงพอ (ทองปาน, 2556)

\section{ความเป็นมาและวัตถุประสงค์ของนโยบาย Smart Farmer}

กระทรวงเกษตรและสหกรณ์จึงได้ริเริ่มโครงการ ``เกษตรกรอัจฉริยะ'' (Smart Farmer: SF) ขึ้นในปี พ.ศ. 2557 เพื่อปฏิรูปภาคการเกษตรไทยสู่ยุคดิจิทัล โดยมุ่งเน้นการพัฒนาองค์ความรู้และส่งเสริมการใช้เทคโนโลยีในการผลิต และมุ่งเน้นให้เกิดการเปลี่ยนผ่านจาก ``เกษตรกรต้นแบบ'' ไปสู่ ``เกษตรกรสมัยใหม่'' ที่เน้นการบริหารจัดการและเทคโนโลยี (Meechoovet \& Siriwato, 2023; เดอะเนชั่นไทยแลนด์, 2566)

ซึ่งโครงการได้สร้างเกษตรกรอัจฉริยะแล้วกว่า 166,900 ราย เกษตรกรอัจฉริยะต้นแบบ 18,143 ราย และกลุ่มเกษตรกรอัจฉริยะ 462 กลุ่ม

และจากการศึกษาเกี่ยวกับผลกระทบของการเข้าร่วมบริการส่งเสริมการเกษตรที่มีต่อประสิทธิภาพของฟาร์ม พบว่า บริการส่งเสริมการเกษตรถือเป็นแนวทางสำคัญในการพัฒนาผลิตภาพทางการเกษตร (Atsan et al., 2009; Egziabher et al., 2011) ผ่านการถ่ายทอดความรู้ด้านเทคโนโลยีและการจัดการฟาร์มให้แก่เกษตรกร
ส่่งผลให้ผู้เข้าร่วมสามารถเพิ่มความสามารถทางการเกษตร นำไปสู่การเพิ่มผลิตภาพและประสิทธิภาพทางธุรกิจอย่างมีนัยสำคัญ \cite{hwang2024}

\subsection{การจำแนกประเภทและคุณสมบัติของ Smart Farmer (KPIs)}

เพื่อให้การพัฒนาเป็นไปอย่างตรงเป้าหมาย กรมส่งเสริมการเกษตรและหน่วยงานที่เกี่ยวข้องได้กำหนดเกณฑ์คุณสมบัติ (Key Performance Indicators: KPIs) ในการจำแนกเกษตรกรออกเป็น 3 กลุ่มหลัก ดังนี้

\textbf{กลุ่มที่ 1: เกษตรกรปราดเปรื่อง (Smart Farmer)}

หมายถึง เกษตรกรทั่วไปที่มีการประกอบอาชีพเกษตรเป็นหลัก และผ่านการประเมินคุณสมบัติพื้นฐาน 6 ประการ (6 Basic Properties) ที่แสดงถึงความเป็นมืออาชีพ ได้แก่

\begin{itemize}
\item ด้านรายได้: มีรายได้จากภาคการเกษตรไม่ต่ำกว่า 180,000 บาท/ครัวเรือน/ปี
\item ด้านองค์ความรู้: มีความรู้จริงในเรื่องที่ทำ สามารถอธิบายกระบวนการผลิตของตนเองได้
\item ด้านการจัดการข้อมูล: มีข้อมูลประกอบการตัดสินใจ มีการจดบันทึกบัญชีฟาร์มและใช้ข้อมูลการตลาดนำการผลิต
\item ด้านการบริหารจัดการ: มีการบริหารจัดการผลผลิตและการตลาดอย่างเป็นระบบ เพื่อลดความเสี่ยงด้านราคา
\item ด้านความปลอดภัย: มีความตระหนักถึงคุณภาพสินค้าและความปลอดภัยของผู้บริโภค (Food Safety) และกระบวนการผลิตที่เป็นมิตรต่อสิ่งแวดล้อม
\item ด้านความรับผิดชอบ: มีความรับผิดชอบต่อสังคมและสิ่งแวดล้อม
\end{itemize}

\textbf{กลุ่มที่ 2: เกษตรกรต้นแบบ (Smart Farmer Model: SFM)}

คือ Smart Farmer ที่ได้รับการคัดเลือกและพัฒนาต่อยอดจนมีศักยภาพสูง สามารถเป็นวิทยากรหรือผู้นำในการถ่ายทอดความรู้ให้แก่เกษตรกรรายอื่นได้ รวมถึงเป็นเจ้าของศูนย์เรียนรู้การเพิ่มประสิทธิภาพการผลิตสินค้าเกษตร (ศพก.)

\textbf{กลุ่มที่ 3: เกษตรกรรุ่นใหม่ (Young Smart Farmer: YSF)}

คือ บุตรหลานเกษตรกรหรือคนรุ่นใหม่ อายุระหว่าง 17--45 ปี ที่มีความมุ่งมั่นตั้งใจทำการเกษตรทดแทนคนรุ่นเก่า มีทักษะในการใช้เทคโนโลยีดิจิทัลและนวัตกรรม (Agri-Tech) ในการบริหารจัดการฟาร์มและการตลาดออนไลน์

จากข้อมูลข้างต้นพบว่า แม้โครงการจะสามารถสร้างเกษตรกรอัจฉริยะได้เป็นจำนวนมาก แต่ข้อมูลจากการสำรวจพบข้อสังเกตสำคัญในเชิงระเบียบวิธีวิจัยว่า เกษตรกรที่เข้าร่วมโครงการมักมีพื้นฐานทักษะเริ่มต้นสูงกว่ากลุ่มที่ไม่เข้าร่วมอย่างมีนัยสำคัญ

สภาวะดังกล่าวนำไปสู่ประเด็นความเอนเอียงในการคัดเลือก (Selection Bias) ซึ่งอาจส่งผลให้การประเมินประสิทธิผลของนโยบายในภาพรวมมีความคลาดเคลื่อน หากไม่มีการควบคุมปัจจัยด้านความพร้อมเบื้องต้นของเกษตรกร \cite{hwang2024}

\subsection{แนวทางการพัฒนาและกระบวนการคัดกรอง}
กระบวนการพัฒนา Smart Farmer ของกรมส่งเสริมการเกษตร ใช้วิธีการประเมินศักยภาพเกษตรกรรายบุคคล (Assessment) เทียบกับเกณฑ์ 6 ข้อข้างต้น เพื่อจัดทำ แผนพัฒนารายบุคคล (Individual Development Plan: IDP) หากเกษตรกรยังขาดคุณสมบัติข้อใด จะได้รับการส่งเสริมและอบรมในด้านนั้นๆ เพิ่มเติม เช่น การอบรมทำบัญชีครัวเรือน การอบรมมาตรฐาน GAP หรือการอบรมเทคโนโลยีการผลิตสมัยใหม่นอกจากนี้ ยังมีการส่งเสริมการรวมกลุ่มในรูปแบบ Smart Group และยกระดับสู่ Smart Enterprise เพื่อสร้างอำนาจต่อรองและเชื่อมโยงเครือข่ายธุรกิจเกษตรครบวงจร

\section{ปัจจัยที่ส่งผลต่อความสามารถและการยอมรับเทคโนโลยีของเกษตรกร}
\subsection{ลักษณะของครัวเรือนและอายุ}
ปัจจัยส่วนใหญ่มีผลเชิงบวกต่อความสามารถทางการเกษตร ยกเว้นอายุของหัวหน้าครัวเรือนและประสบการณ์ทำฟาร์ม \cite{abdallah2016,anang2020}
โดยผู้สูงอายุมักมีข้อจำกัดในการเรียนรู้และยอมรับเทคโนโลยีใหม่ เนื่องจากความกังวลด้านความยุ่งยากในการใช้งาน (Perceived Ease of Use) ส่งผลให้มีแนวโน้มเข้าร่วมโครงการ SF น้อยกว่า แต่สามารถบรรเทาได้ด้วยระบบการเรียนรู้ทางไกล หรือ e-learning (Duang-ek-Anong, 2019)
\subsection{ผลลองประสบการณ์ทำฟาร์ม}
งานวิจัยแสดงผลที่แตกต่างกัน บางการศึกษาพบว่าเกษตรกรที่มีประสบการณ์มากมีแนวโน้มเข้าร่วมบริการส่งเสริมการเกษตรและพัฒนาทักษะได้ดีขึ้น (Eticha, 2021; Anang \& Ayambila, 2020; Goshu, 2019) ในขณะที่บางการศึกษาพบว่าประสบการณ์มีผลเชิงลบ (Abdallah \& Abdul-Rahaman, 2016) หรือไม่มีผลต่อทักษะทางการเกษตร (Worku, 2022)

\subsection{ลักษณะของฟาร์ม}
ปัจจัยส่วนใหญ่มีผลเชิงบวกต่อระดับการเข้าร่วมบริการส่งเสริมการเกษตร ยกเว้นการทำกิจกรรมนอกฟาร์ม โดยมีผลการศึกษาที่ขัดแย้งกัน ระหว่างผลกระทบเชิงลบ (Abegunde et al., 2019) และผลกระทบเชิงบวก (Nahayo et al., 2017)

\subsection{สภาพแวดล้อมและสภาวะตลาด}
เกษตรกรที่เผชิญความเสี่ยงสูงในการผลิตและขาดแคลนทรัพยากรมักมีความสามารถทางการเกษตรต่ำ โดย Nahayo et al. (2017) พบว่าการขาดแคลนระบบชลประทาน เครื่องจักรกล บริการส่งเสริม และปัจจัยการผลิต รวมถึงเทคโนโลยี ส่งผลเชิงลบต่อการเข้าร่วมโครงการ
\subsection{ความสัมพันธ์กับความเสี่ยง}
Ghadim et al. (2005) ชี้ให้เห็นว่า ขีดความสามารถของเกษตรกรในการแสวงหาข้อมูลและกระบวนการเรียนรู้เทคโนโลยีใหม่ มีความสัมพันธ์อย่างมีนัยสำคัญกับทัศนคติที่มีต่อความเสี่ยง (Risk Preference) ในการทำการเกษตร
อย่างไรก็ตาม ในบริบทของการดำเนินนโยบายส่งเสริม SF นั้น เกษตรกรมีอิสระในการเลือกเข้าร่วมโครงการตามความสมัครใจ (Self-selection Bias) ซึ่ง Hwang et al. (2024) ระบุว่าปัจจัยดังกล่าวอาจก่อให้เกิดปัญหาความแตกต่างระหว่างกลุ่มเป้าหมาย (Heterogeneity) ตั้งแต่จุดเริ่มต้นในการเรียนรู้
ซึ่งนำไปสู่การประสบความสำเร็จในการประยุกต์ใช้เทคโนโลยีเข้ามาช่วยในการทำการเกษตรในระยะยาว นอกจากนี้ การประเมินผลสัมฤทธิ์ของนโยบาย SF ไม่สามารถอาศัยเพียงการเปรียบเทียบค่าเฉลี่ยก่อนและหลังการเข้าร่วมโครงการ (Simple Pre-Post Comparison) ได้โดยตรง เนื่องจากความซับซ้อนของพฤติกรรมมนุษย์และอิทธิพลจากปัจจัยภายนอกที่ควบคุมไม่ได้ (Unobserved Factors)

\section{บทบาทของ PSM และ DID ในการประเมินนโยบายเชิงเศรษฐศาสตร์}
สืบเนื่องจากหัวข้อ 3 ที่เกี่ยวข้องกับการประเมินผลสัมฤทธิ์ของนโยบาย Smart Farmer ของภาครัฐ ซึ่งโดยพื้นฐานแล้วไม่สามารถวัดเพียงการเปรียบเทียบค่าเฉลี่ยก่อนและหลังโครงการได้ เนื่องจากความซับซ้อนของพฤติกรรมมนุษย์และปัจจัยภายนอก นักเศรษฐศาสตร์จึงนิยมใช้แนวทาง Quasi-experimental Design ผ่านเทคนิค PSM และ DID โดยมีเป้าหมายหลักดังนี้:

\subsection{การจัดการปัญหา Selection Bias ด้วย PSM}
เป้าหมายสำคัญของ Propensity Score Matching (PSM) คือการสร้าง ``กลุ่มเปรียบเทียบที่เสมือนจริง'' (Counterfactual) เนื่องจากในความเป็นจริง เกษตรกรไม่ได้ถูกสุ่มเข้าโครงการ แต่เลือกเข้าร่วมด้วยตนเอง (Self-selection)
ซึ่ง Hwang et al. (2024) ระบุว่าเกษตรกรกลุ่มนี้มักมี ``ต้นทุนทางสังคม'' หรือความพร้อมสูงกว่าปกติอยู่แล้ว ซึ่งสอดคล้องกับงานวิจัยของ Olounlade et al. (2020) ที่ชี้ให้เห็นว่าปัจจัยด้านลักษณะทางประชากรและทรัพยากรที่แตกต่างกันทำให้เกิดความลำเอียงในการคัดเลือก (Selection Bias)
ซึ่งหากไม่ได้รับการควบคุม จะส่งผลให้การประเมินค่าผลกระทบของนโยบายคลาดเคลื่อนจากความเป็นจริง
ด้วยเหตุนี้ การตรวจสอบผลผลิตที่เพิ่มขึ้น (Outcome) ว่าเกิดจาก ``นโยบาย'' ไม่ใช่เกิดจาก ``คุณลักษณะเดิม'' ของเกษตรกรที่มีความเก่งหรือมีทุนหนาอยู่แล้ว
การทำ PSM จึงเป็นการปรับพื้นฐาน (Baseline) ของกลุ่มตัวอย่างให้เท่ากันก่อนการเปรียบเทียบ โดยใช้ค่า Propensity Score เป็นตัวกลางในการจับคู่เกษตรกรที่มีลักษณะใกล้เคียงกันที่สุด (Adam et al., 2022)

\subsection{การแยกแยะผลกระทบแท้จริงด้วย DID}
เป้าหมายของ Difference-in-Differences (DID) คือการขจัดอิทธิพลของปัจจัยภายนอกที่เปลี่ยนแปลงตามเวลา (Time-varying factors) เช่น สภาพอากาศในแต่ละปี หรือภาวะเศรษฐกิจโลกที่กระทบเกษตรกรทุกคนเหมือนกัน โดยงานวิจัยของ Li และ Yang (2024) ได้ใช้แนวทาง PSM-DID เพื่อประเมินผลกระทบของการเข้าร่วมสหกรณ์ต่อรายได้เกษตรกรในจีน ซึ่งพิสูจน์ให้เห็นว่าวิธีนี้สามารถแยกแยะผลกระทบจากนโยบายออกจากปัจจัยภายนอกและปัจจัยที่ไม่เปลี่ยนแปลงตามเวลา (Unobserved Time-invariant Factors) ได้อย่างมีประสิทธิภาพ โดย DID ถูกใช้เพื่อระบุผลกระทบสุทธิ (Net Impact) โดยการเปรียบเทียบการเปลี่ยนแปลงของกลุ่มที่ได้รับนโยบายเทียบกับกลุ่มที่ไม่ได้รับนโยบายในช่วงเวลาเดียวกัน ซึ่งช่วยตอบโจทย์ว่านโยบาย Smart Farmer สามารถสร้างความแตกต่าง (Value-added) ได้จริงหรือไม่ เมื่อเทียบกับกรณีที่ไม่มีนโยบายนี้ อย่างไรก็ตาม ในบางกรณี เช่น การศึกษาของ Adam et al. (2022) ในอินโดนีเซีย แม้จะใช้ PSM-DID แล้ว แต่ผลลัพธ์กลับพบว่านโยบายอาจไม่ส่งผลเชิงบวกต่อผลิตภาพเสมอไป หากปัจจัยเชิงโครงสร้างอื่นๆ ไม่ได้รับการแก้ไข

โดยสรุปแล้ว PSM และ DID เป็นเครื่องมือที่มีประสิทธิภาพในการหาในการหา Causal Inference (ความเป็นเหตุเป็นผล) เพื่อตอบวัตถุประสงค์ว่านโยบายส่งผลต่อผลิตภาพในภาพรวมหรือไม่ แต่เมื่อผลการวิเคราะห์พบว่าไม่มีนัยสำคัญทางสถิติ (Non-significant) สิ่งนี้กลับสะท้อนถึงประเด็นทางเศรษฐศาสตร์ที่ลึกซึ้งขึ้น นั่นคือ 'ความล้มเหลวของค่าเฉลี่ย' (The Flaw of Averages)" โดยในบริบทของ Smart Farmer ความสำเร็จอาจไม่ได้กระจายตัวแบบสม่ำเสมอ แต่อาจกระจุกตัวอยู่ในเกษตรกรบางกลุ่มที่มีเงื่อนไขเฉพาะ (Heterogeneous Treatment Effects) ด้วยเหตุนี้ งานวิจัยจึงจำเป็นต้องข้ามพ้นขอบเขตของแบบจำลองเชิงเส้น ไปสู่การใช้ Machine Learning (ML) เพื่อทำการ Classification เป้าหมายเปลี่ยนจากการถามว่า 'นโยบายหรือหลักสูตรที่สร้างขึ้นมานั้นได้ผลหรือไม่' (Average Impact) ไปสู่การระบุว่า 'นโยบายหรือหลักสูตรนี้เหมาะสมกับใครและภายใต้เงื่อนไขใดบ้าง' (Pattern Identification) ซึ่งจะเป็นการเติมเต็มวัตถุประสงค์ของงานวิจัยในการเสนอแนะนโยบายแบบตรงจุด (Targeted Policy) มากขึ้น

\subsection{การเปลี่ยนผ่านสู่ Machine Learning เพื่อวิเคราะห์ผลกระทบที่ไม่สม่ำเสมอ (Heterogeneity)}
โดยสรุปแล้ว PSM และ DID เป็นเครื่องมือที่มีประสิทธิภาพในการหา Causal Inference (ความเป็นเหตุเป็นผล) เพื่อตอบวัตถุประสงค์ว่านโยบายส่งผลต่อผลิตภาพในภาพรวมหรือไม่ แต่เมื่อผลการวิเคราะห์พบว่าไม่มีนัยสำคัญทางสถิติ (Non-significant) สิ่งนี้สะท้อนให้เห็นถึง นั่นคือ "ความล้มเหลวของค่าเฉลี่ย" (The Flaw of Averages) ซึ่งอาจจะเกิดได้จากการกระจายตัวที่ไม่สม่ำเสมอของข้อมูล โดยมีการกระจุกตัวในเกษตรกรบางกลุ่มที่มีเงื่อนไขเฉพาะ (Heterogeneous Treatment Effects)  ดังเช่นในงานวิจัยของ Li and Yang (2024) พบว่าผลกระทบของนโยบายจะเด่นชัดกว่าในกลุ่มเกษตรกรที่มีรายได้น้อยถึงปานกลาง พบว่าผลกระทบของนโยบายจะเด่นชัดกว่าในกลุ่มเกษตรกรที่มีรายได้น้อยถึงปานกลาง ด้วยเหตุนี้ งานวิจัยจึงจำเป็นต้องข้ามพ้นขอบเขตของแบบจำลองเชิงเส้น ไปสู่การใช้ Machine Learning (ML) เพื่อทำการ Classification เป้าหมายเปลี่ยนจากการถามว่า "นโยบายที่ภาครัฐดำเนินการส่งเสริมอยู่นั้น “ได้ผลหรือไม่" (Average Impact) ไปสู่คำถามว่า "นโยบายที่ดำเนินไปอยู่นี้มีผลกับใคร และภายใต้เงื่อนไขใด" (Pattern Identification) โดยอาศัยการระบุตัวแปรสำคัญ (Feature Importance) ด้วย ML มาจำแนกรูปแบบพฤติกรรมและความพร้อมทางเทคโนโลยีที่นำไปสู่ผลสัมฤทธิ์ ซึ่งจะเป็นการเติมเต็มวัตถุประสงค์ของงานวิจัยในการเสนอแนะนโยบายแบบตรงจุด (Targeted Policy) มากขึ้น



\section{เทคโนโลยีการเรียนรู้ของเครื่อง (Machine Learning)}
การเรียนรู้ของเครื่อง (Machine Learning) เป็นประเภทหนึ่งของปัญญาประดิษฐ์ (AI) ที่ช่วยให้คอมพิวเตอร์สามารถเรียนรู้และพัฒนาจากประสบการณ์ได้โดยไม่ต้องเขียนโปรแกรมอย่างชัดเจน โดยใช้อัลกอริทึมและโมเดลทางสถิติในการวิเคราะห์และตีความข้อมูล เพื่อทำนายหรือตัดสินใจจากการวิเคราะห์นั้นการเรียนรู้ของเครื่องมี 3 ประเภทหลัก ได้แก่:

\begin{table}[htbp]
\centering
\small
\setlength{\tabcolsep}{4pt}
\renewcommand{\arraystretch}{1.2}

\caption{ลักษณะภาพรวมของประเภทการเรียนรู้ของเครื่อง}
\label{tab:ml_types}

\begin{tabular}{p{3cm} p{6cm} p{6cm}}
\hline
\textbf{ประเภท} & \textbf{Advantages} & \textbf{Disadvantages} \\
\hline

Supervised Learning &
\begin{itemize}
\item สามารถทำนายได้แม่นยำสูงโดยใช้ข้อมูลที่มีการติดป้าย
\item ง่ายต่อการนำไปใช้และตีความ
\item สามารถจัดการได้ทั้งงานการถดถอยและการจำแนกประเภท
\end{itemize}
&
\begin{itemize}
\item ต้องใช้ข้อมูลที่มีการติดป้ายในการฝึกสอน ซึ่งใช้เวลาและมีค่าใช้จ่ายสูง
\item อาจเกิดการเรียนรู้มากเกินไป (Overfit) หากโมเดลซับซ้อนเกินไป
\item อาจทำงานได้ไม่ดีหากข้อมูลมีสัญญาณรบกวนหรือค่าผิดปกติ
\end{itemize}
\\

Unsupervised Learning &
\begin{itemize}
\item สามารถค้นหารูปแบบและความสัมพันธ์ที่ซ่อนอยู่ในข้อมูล
\item ไม่ต้องใช้ข้อมูลที่มีการติดป้ายในการฝึกสอน
\item สามารถจัดการกับชุดข้อมูลขนาดใหญ่ได้
\end{itemize}
&
\begin{itemize}
\item ผลลัพธ์อาจยากต่อการตีความและตรวจสอบ
\item อาจทำงานได้ไม่ดีกับชุดข้อมูลขนาดเล็กหรือข้อมูลที่มีมิติสูง
\item อ่อนไหวต่อการกำหนดค่าเริ่มต้นและพารามิเตอร์
\end{itemize}
\\

Reinforcement Learning &
\begin{itemize}
\item สามารถเรียนรู้นโยบายที่เหมาะสมสำหรับงานที่ซับซ้อน
\item สามารถจัดการกับสภาพแวดล้อมที่มีการเปลี่ยนแปลง
\item สามารถเรียนรู้จากการลองผิดลองถูกโดยไม่ต้องมีการกำกับดูแลอย่างชัดเจน
\end{itemize}
&
\begin{itemize}
\item ใช้ทรัพยากรคอมพิวเตอร์สูงและใช้เวลามาก
\item ต้องมีฟังก์ชันรางวัลที่กำหนดไว้อย่างชัดเจน
\item อาจมีปัญหาในการสมดุลระหว่างการสำรวจและการใช้ประโยชน์
\end{itemize}
\\

\hline
\end{tabular}
\end{table}

สำหรับงานวิจัยนี้ เลือกใช้การเรียนรู้แบบมีผู้สอน เนื่องจากมีชุดข้อมูลเกี่ยวกับเกษตรกรอัจฉริยะในประเทศไทยและทราบถึงทักษะที่จำเป็นในการเป็นเกษตรกรอัจฉริยะ ซึ่งเหมาะสมกับการใช้ข้อมูลเฉพาะจากเกษตรกร โดยในเบื้องต้น พบว่า การประยุกต์ใช้ ML เพื่อพัฒนาทักษะด้านการพัฒนาอาชีพนั้น เป็นเรื่องที่ได้รับการนิยมในหลากหลายสาขาวิชา อาทิ สาขาการศึกษา ที่ได้มีการนำเสนอการใช้ Supervised Learning ในการวิเคราะห์ความสัมพันธ์ระหว่างคุณลักษณะของผู้เรียนกับผลสัมฤทธิ์ทางการเรียน นำไปสู่การออกแบบหลักสูตรที่ตอบสนองความต้องการเฉพาะบุคคล ขณะที่ Premanode (2556) แสดงให้เห็นว่าการใช้ Machine Learning ในการติดตามและประเมินผลการเรียนรู้แบบ Real-time ช่วยให้สามารถปรับปรุงเนื้อหาและวิธีการสอนได้อย่างทันท่วงที รวมไปถึงด้านการเกษตร Charoensukmongkol (2558) พบว่า การใช้ Machine Learning ช่วยในการวิเคราะห์รูปแบบการเรียนรู้ของเกษตรกรและปัจจัยที่ส่งผลต่อการยอมรับเทคโนโลยีใหม่ ซึ่งนำไปสู่การพัฒนาหลักสูตรฝึกอบรมที่มีประสิทธิภาพมากขึ้น สอดคล้องกับงานของ Jinjarak (2550) ที่ชี้ให้เห็นว่าการใช้ระบบอัจฉริยะในการจำแนกกลุ่มเกษตรกรตามความต้องการและข้อจำกัด ช่วยเพิ่มประสิทธิผลของโครงการพัฒนาทักษะ อย่างไรก็ตาม ก็มีการแสดงถึงข้อจำกัดการใช้ ML ในบางกลุ่มข้อมูลที่เกี่ยวข้องกับการพัฒนาอาชีพ โดย Ruete (2557) ชี้ให้เห็นข้อจำกัดของการใช้ Machine Learning ในการพัฒนาทักษะอาชีพ โดยเฉพาะในกลุ่มผู้เรียนที่มีข้อจำกัดด้านการเข้าถึงเทคโนโลยีและทักษะดิจิทัล จึงจำเป็นต้องพัฒนาระบบที่เข้าถึงง่ายและมีการสนับสนุนที่เหมาะสม
จากการทบทวนวรรณกรรมข้างต้นสามารถสรุปได้ว่า Machine Learning เป็นเครื่องมือที่มีศักยภาพในการวิเคราะห์ข้อมูลที่มีความซับซ้อนสูง และสามารถอธิบายความแตกต่างของพฤติกรรมหรือผลลัพธ์ในระดับรายบุคคลได้ ซึ่งเหมาะสมอย่างยิ่งสำหรับบริบทภาคการเกษตร ที่มีความหลากหลายของเกษตรกรทั้งด้านศักยภาพและทรัพยากร